\documentclass[aspectratio=169]{beamer}

%----------------- GÓI LỆNH CƠ BẢN -----------------
\usepackage[utf8]{inputenc}
\usepackage{lmodern}
\usepackage[T5]{fontenc}
\usepackage[vietnamese]{babel}
\usepackage{ragged2e}
\usepackage{tikz}
\usetikzlibrary{calc, arrows.meta, positioning, shapes, shapes.geometric}
\usepackage{listings}
\usepackage{xcolor}

% Cấu hình hiển thị Code C++
\lstset{
    language=C++,
    basicstyle=\ttfamily\tiny,
    keywordstyle=\color{blue}\bfseries,
    commentstyle=\color{green!60!black},
    stringstyle=\color{orange},
    numbers=left,
    numberstyle=\tiny\color{gray},
    stepnumber=1,
    frame=single,
    breaklines=true,
    backgroundcolor=\color{gray!5}
}

\tikzset{
    state/.style={circle, draw, thick, minimum size=0.7cm, fill=#1!20},
    state/.default=blue,
    node distance=1.5cm
}

\usetheme{Madrid}
\usecolortheme{default}

%----------------- THÔNG TIN TIÊU ĐỀ -----------------
\title[Lý thuyết Đồ thị & Thuật toán]{LÝ THUYẾT ĐỒ THỊ VÀ THUẬT TOÁN}
\subtitle{Cơ sở dữ liệu - Giải thuật - Cài đặt C++}

\date[\today]{\today}

\begin{document}

% 1. Trang tiêu đề
\begin{frame}
    \titlepage
\end{frame}

% 2. Mục lục
\begin{frame}{Nội dung báo cáo}
    \tableofcontents
\end{frame}

% --- PHẦN 1: GIỚI THIỆU ---
\section{Giới thiệu & Định nghĩa}
% 3. Khái niệm
\begin{frame}{1.1 Khái niệm về Đồ thị}
    \begin{itemize}
        \item Đồ thị là cấu trúc dữ liệu gồm các nút và các kết nối giữa chúng.
        \item \textbf{Bài toán thực tế:}
        \begin{itemize}
            \item Mạng internet: Các trang web (đỉnh) và liên kết (cạnh).
            \item Bản đồ: Các thành phố (đỉnh) và đường đi (cạnh).
        \end{itemize}
    \end{itemize}
    
\end{frame}

% 4. Định nghĩa toán học
\begin{frame}{1.2 Định nghĩa toán học}
    \begin{block}{Đồ thị $G = (V, E)$}
        \begin{itemize}
            \item $V = \{v_1, v_2, \dots, v_n\}$ là tập các \textbf{đỉnh} (Vertices).
            \item $E = \{e_1, e_2, \dots, e_m\}$ là tập các \textbf{cạnh} (Edges). Mỗi cạnh $e \in E$ nối một cặp đỉnh $(u, v)$.
        \end{itemize}
    \end{block}
    \centering
    \begin{tikzpicture}
        \node[state] (A) {A};
        \node[state, right of=A] (B) {B};
        \node[state, below of=A] (C) {C};
        \draw[thick] (A) -- (B);
        \draw[thick] (B) -- (C);
        \draw[thick] (A) -- (C);
    \end{tikzpicture}
    \\ \small Ví dụ: $V = \{A, B, C\}, E = \{(A,B), (B,C), (A,C)\}$
\end{frame}

% --- PHẦN 2: PHÂN LOẠI ---
\section{Phân loại đồ thị}
% 5. Đồ thị vô hướng
\begin{frame}{2.1 Đồ thị vô hướng (Undirected Graph)}
    \begin{columns}
        \begin{column}{0.6\textwidth}
            \begin{itemize}
                \item Cạnh nối hai đỉnh không phân biệt thứ tự.
                \item Cạnh $(u, v)$ cũng là cạnh $(v, u)$.
                \item Biểu diễn quan hệ hai chiều (ví dụ: Bạn bè trên Facebook).
            \end{itemize}
        \end{column}
        \begin{column}{0.4\textwidth}
            \centering
            \begin{tikzpicture}
                \node[state] (1) {1};
                \node[state, right of=1] (2) {2};
                \draw[thick] (1) -- (2);
            \end{tikzpicture}
        \end{column}
    \end{columns}
\end{frame}

% 6. Đồ thị có hướng
\begin{frame}{2.2 Đồ thị có hướng (Directed Graph/Digraph)}
    \begin{columns}
        \begin{column}{0.6\textwidth}
            \begin{itemize}
                \item Cạnh có hướng xác định từ đỉnh nguồn đến đỉnh đích.
                \item Ký hiệu cung: $u \to v$.
                \item Ví dụ: Lượt theo dõi trên Instagram, đường một chiều.
            \end{itemize}
        \end{column}
        \begin{column}{0.4\textwidth}
            \centering
            \begin{tikzpicture}
                \node[state] (1) {1};
                \node[state, right of=1] (2) {2};
                \draw[thick, ->, >={Stealth}] (1) -- (2);
            \end{tikzpicture}
        \end{column}
    \end{columns}
    
\end{frame}

% 7. Đồ thị có trọng số
\begin{frame}{2.3 Đồ thị có trọng số (Weighted Graph)}
    \begin{itemize}
        \item Mỗi cạnh gắn liền với một giá trị số (trọng số).
        \item Trọng số có thể là: khoảng cách, chi phí, thời gian,...
    \end{itemize}
    \centering
    \begin{tikzpicture}
        \node[state] (1) {A};
        \node[state, right=2cm of 1] (2) {B};
        \draw[thick] (1) -- node[above] {50km} (2);
    \end{tikzpicture}
\end{frame}

% --- PHẦN 3: TÍNH CHẤT ---
\section{Tính chất cơ bản}
% 8. Bậc của đỉnh
\begin{frame}{3.1 Bậc của đỉnh}
    \begin{itemize}
        \item \textbf{Vô hướng:} $deg(v)$ là số cạnh kề với đỉnh $v$.
        \item \textbf{Có hướng:} 
        \begin{itemize}
            \item Bậc vào ($deg^-(v)$): số cung đi vào.
            \item Bậc ra ($deg^+(v)$): số cung đi ra.
        \end{itemize}
    \end{itemize}
    \begin{exampleblock}{Định lý bắt tay}
        Tổng bậc của tất cả các đỉnh bằng 2 lần số cạnh: $\sum deg(v) = 2|E|$.
    \end{exampleblock}
\end{frame}

% --- PHẦN 4: TỔ CHỨC DỮ LIỆU ---
\section{Tổ chức dữ liệu}
% 9. Ma trận kề
\begin{frame}{4.1 Ma trận kề (Adjacency Matrix)}
    \begin{columns}
        \begin{column}{0.5\textwidth}
            \begin{itemize}
                \item Ma trận $A_{n \times n}$.
                \item $A[i][j] = 1$ nếu có cạnh, ngược lại $0$.
                \item \textbf{Ưu:} Kiểm tra cạnh nhanh $O(1)$.
                \item \textbf{Nhược:} Tốn bộ nhớ $O(V^2)$.
            \end{itemize}
        \end{column}
        \begin{column}{0.5\textwidth}
            $A = \begin{pmatrix} 0 & 1 & 1 \\ 1 & 0 & 0 \\ 1 & 0 & 0 \end{pmatrix}$
        \end{column}
    \end{columns}
\end{frame}

% 10. Danh sách kề
\begin{frame}{4.2 Danh sách kề (Adjacency List)}
    \begin{itemize}
        \item Mỗi đỉnh có một danh sách các đỉnh kề nó.
        \item Sử dụng \texttt{vector<int> adj[N]} trong C++.
        \item \textbf{Ưu điểm:} Tiết kiệm bộ nhớ, duyệt các đỉnh kề nhanh.
    \end{itemize}
    
\end{frame}

% 11. Ví dụ Code cài đặt danh sách kề
\begin{frame}[fragile]{4.3 Ví dụ Code: Khai báo đồ thị}
\begin{lstlisting}
#include <iostream>
#include <vector>
using namespace std;

const int MAXN = 1000;
vector<int> adj[MAXN]; // Mang cac vector de luu danh sach ke

void addEdge(int u, int v) {
    adj[u].push_back(v);
    adj[v].push_back(u); // Bo dong nay neu la do thi co huong
}

int main() {
    addEdge(1, 2);
    addEdge(1, 3);
    return 0;
}
\end{lstlisting}
\end{frame}

% --- PHẦN 5: DFS ---
\section{Thuật toán DFS}
% 12. Ý tưởng DFS
\begin{frame}{5.1 Tìm kiếm theo chiều sâu (DFS)}
    \begin{itemize}
        \item \textbf{Nguyên lý:} Đi sâu nhất có thể trên một nhánh, khi không đi được nữa thì quay lui (Backtracking).
        \item \textbf{Cơ chế:} Sử dụng Ngăn xếp (Stack) hoặc Đệ quy.
    \end{itemize}
    
\end{frame}

% 13. Code DFS
\begin{frame}[fragile]{5.2 Cài đặt DFS bằng đệ quy}
\begin{lstlisting}
bool visited[MAXN];

void DFS(int u) {
    cout << u << " "; // Tham dinh u
    visited[u] = true;
    
    for (int v : adj[u]) {
        if (!visited[v]) {
            DFS(v); // Goi de quy den dinh ke chua tham
        }
    }
}
\end{lstlisting}
\end{frame}

% 14. Ví dụ minh họa DFS
\begin{frame}{5.3 Ví dụ minh họa DFS}
    \centering
    \begin{tikzpicture}
        \node[state] (1) {1};
        \node[state, below left of=1] (2) {2};
        \node[state, below right of=1] (3) {3};
        \node[state, below of=2] (4) {4};
        \draw[thick] (1) -- (2) -- (4);
        \draw[thick] (1) -- (3);
    \end{tikzpicture}
    \\ \textbf{Thứ tự duyệt:} 1 $\to$ 2 $\to$ 4 $\to$ 3
\end{frame}

% --- PHẦN 6: BFS ---
\section{Thuật toán BFS}
% 15. Ý tưởng BFS
\begin{frame}{6.1 Tìm kiếm theo chiều rộng (BFS)}
    \begin{itemize}
        \item \textbf{Nguyên lý:} Thăm các đỉnh theo từng lớp (Layer). Thăm hết các đỉnh cách gốc $k$ cạnh trước khi thăm các đỉnh cách gốc $k+1$.
        \item \textbf{Cơ chế:} Sử dụng Hàng đợi (Queue).
    \end{itemize}
    
\end{frame}

% 16. Code BFS
\begin{frame}[fragile]{6.2 Cài đặt BFS}
\begin{lstlisting}
#include <queue>

void BFS(int startNode) {
    queue<int> q;
    q.push(startNode);
    visited[startNode] = true;

    while (!q.empty()) {
        int u = q.front(); q.pop();
        cout << u << " ";

        for (int v : adj[u]) {
            if (!visited[v]) {
                visited[v] = true;
                q.push(v);
            }
        }
    }
}
\end{lstlisting}
\end{frame}

% 17. Ví dụ minh họa BFS
\begin{frame}{6.3 Ví dụ minh họa BFS}
    \centering
    \begin{tikzpicture}
        \node[state] (1) {1};
        \node[state, below left of=1] (2) {2};
        \node[state, below right of=1] (3) {3};
        \node[state, below of=2] (4) {4};
        \draw[thick] (1) -- (2) -- (4);
        \draw[thick] (1) -- (3);
    \end{tikzpicture}
    \\ \textbf{Thứ tự duyệt:} 1 $\to$ 2 $\to$ 3 $\to$ 4 (vì 2 và 3 cùng mức)
\end{frame}

% --- PHẦN 7: SO SÁNH & ỨNG DỤNG ---
\section{So sánh & Ứng dụng}
% 18. So sánh DFS và BFS
\begin{frame}{7.1 So sánh DFS và BFS}
    \begin{table}
        \begin{tabular}{|l|l|l|}
            \hline
            \textbf{Đặc điểm} & \textbf{DFS} & \textbf{BFS} \\ \hline
            Cấu trúc & Stack / Đệ quy & Queue \\ \hline
            Đường đi ngắn nhất & Không (với đồ thị ko trọng số) & \textbf{Có} \\ \hline
            Bộ nhớ & Ít hơn (theo chiều sâu) & Nhiều hơn (theo chiều rộng) \\ \hline
            Ứng dụng & Tìm chu trình, giải đố & Tìm đường ngắn nhất, GPS \\ \hline
        \end{tabular}
    \end{table}
\end{frame}

% 19. Ứng dụng: Tìm đường đi
\begin{frame}[fragile]{7.2 Ứng dụng: Kiểm tra tính liên thông}
    \begin{itemize}
        \item Đồ thị liên thông nếu từ một đỉnh bất kỳ có thể đi tới tất cả các đỉnh khác.
        \item \textbf{Cách làm:} Gọi DFS(1), sau đó kiểm tra mảng \texttt{visited}.
    \end{itemize}
\begin{lstlisting}
int count = 0;
DFS(1);
for(int i=1; i<=n; i++) if(visited[i]) count++;
if(count == n) cout << "Do thi lien thong";
\end{lstlisting}
\end{frame}

% 20. Độ phức tạp
\begin{frame}{7.3 Độ phức tạp thuật toán}
    Đối với cả DFS và BFS sử dụng danh sách kề:
    \begin{itemize}
        \item \textbf{Thời gian:} $O(V + E)$
        \begin{itemize}
            \item $V$: số đỉnh.
            \item $E$: số cạnh.
        \end{itemize}
        \item \textbf{Không gian:} $O(V)$ để lưu trữ mảng đánh dấu và hàng đợi/ngăn xếp.
    \end{itemize}
\end{frame}

% 21. Kết luận
\begin{frame}{Tổng kết}
    \begin{itemize}
        \item Đồ thị là công cụ mạnh mẽ để mô hình hóa thế giới thực.
        \item Lựa chọn cấu trúc dữ liệu (Ma trận hay Danh sách kề) quyết định hiệu suất bài toán.
        \item DFS và BFS là nền tảng cho các thuật toán nâng cao (Dijkstra, Prim, Kruskal).
    \end{itemize}
\end{frame}

% 22. Cảm ơn
\begin{frame}
    \begin{center}
        \textcolor[rgb]{0.50,0.00,0.50}{\Huge \textbf{Trân trọng cảm ơn!}}
        \\ \vspace{1cm}
        \large Bạn có câu hỏi nào không?
    \end{center}
\end{frame}

\end{document}